% =============================================================================
% CAPÍTULO V — CONCLUSIONES
% Proyecto: Kiddy Quiz — UTEQ
% Autor: Jeandavid Cabrera Guerra
% =============================================================================

% ─── PÁRRAFO DE APERTURA DEL CAPÍTULO ────────────────────────────────────────
\noindent
El presente capítulo sintetiza los hallazgos derivados del desarrollo y la
validación empírica de la aplicación web \textit{Kiddy Quiz}. Se presentan,
en primer lugar, las conclusiones del proceso de investigación y desarrollo;
a continuación, los resultados se discuten en contraste con la literatura
previa y se exponen las contribuciones del estudio; seguidamente se reconocen
las limitaciones del trabajo; y, finalmente, se delinean las líneas de
investigación y desarrollo futuro.

% =============================================================================
\section{Conclusiones}
\label{sec:conclusiones}
% =============================================================================

La aplicación de entrevistas semiestructuradas a cinco docentes de educación
básica elemental y a dos psicopedagogos permitió consolidar un cuerpo de
requisitos funcionales y pedagógicos que sustentó la totalidad de la fase de
diseño. De este proceso emergieron veintisiete historias de usuario y seis
requerimientos no funcionales, los cuales reflejan con precisión tanto las
barreras cognitivas reportadas por los profesionales ---particularmente la
lectoescritura emergente y la ansiedad matemática--- como las limitaciones
operativas que enfrentan los docentes al gestionar evaluaciones diferenciadas
en aulas numerosas. El proceso de levantamiento de requisitos se ancló
deliberadamente en la realidad de la educación básica elemental ecuatoriana y
en el Currículo Priorizado con Énfasis en Competencias del Ministerio de
Educación~\cite{mineduc2021curriculo}, garantizando que la solución respondiera
a necesidades verificadas y no a supuestos teóricos.

A partir de dichos requisitos, se diseñó y desarrolló un sistema funcional
construido sobre una arquitectura cliente-servidor ---Angular, NestJS y
PostgreSQL--- y desplegado en producción. La plataforma implementa siete tipos
de preguntas integrables en una misma evaluación, incorpora pictogramas de la
base ARASAAC como soporte cognitivo visual, integra retroalimentación
motivacional generada por el modelo Gemini~2.5~Flash y articula mecanismos de
gamificación basados en niveles e insignias. La adopción de la metodología
Extreme Programming~\cite{beck2004xp,dyba2008agile} permitió incorporar los
ajustes derivados de la retroalimentación continua de los expertos a lo largo
de las treinta y dos semanas del cronograma, ratificando la pertinencia del
enfoque ágil para el desarrollo de software educativo con usuarios de
necesidades específicas.

La validación del sistema, efectuada mediante la Escala de Usabilidad del
Sistema (SUS)~\cite{brooke1996sus} y el Modelo de Aceptación Tecnológica
(TAM)~\cite{davis1989tam} con una muestra de quince estudiantes de educación
básica elemental y cinco docentes en ejercicio, arrojó resultados que superan
los umbrales de referencia en ambas dimensiones. La puntuación media SUS
alcanzó $87{,}7$ puntos en el grupo de estudiantes ($\sigma=7{,}8$;
mín.~$=80{,}0$) y $81{,}0$ en el grupo de docentes ($\sigma=9{,}4$), ubicando
al sistema en la categoría \textbf{``Excelente''} según los rangos
interpretativos de Vlachogianni y Tselios~\cite{vlachogianni2022sus}. En
cuanto a la aceptación tecnológica, la dimensión de Intención de Uso registró
una media de $4{,}93/5{,}0$ entre los estudiantes (mediana~$=5{,}0$;
mín.~$=4{,}6$), lo que indica una disposición prácticamente unánime hacia el
uso continuo de la plataforma. La brecha más pronunciada entre perfiles se
observó en la Facilidad de Uso Percibida, donde los estudiantes ($M=4{,}78$)
superaron a los docentes ($M=3{,}95$), confirmando que el diseño centrado en el
usuario infantil ---con pictogramas ARASAAC y tipografías de gran tamaño---
redujo eficazmente las barreras cognitivas de los niños sin comprometer la
funcionalidad requerida por el docente.

En definitiva, el conjunto de la evidencia recopilada demuestra que es posible
articular, en una única plataforma desplegable, los componentes que la
literatura reciente identifica como prioritarios para la evaluación por
competencias en la infancia: accesibilidad cognitiva, diversidad evaluativa,
retroalimentación inteligente y motivación intrínseca. \textit{Kiddy Quiz}
aborda simultáneamente las cuatro brechas identificadas en el estado del arte:
el nivel educativo objetivo, la diversidad y creación interna de evaluaciones,
la incorporación de retroalimentación generativa con inteligencia artificial y
la articulación de mecanismos de gamificación con accesibilidad cognitiva
mediante pictogramas, acreditando en términos funcionales, técnicos y empíricos
que la tecnología puede constituirse en un puente real hacia la equidad
evaluativa en la educación básica elemental ecuatoriana.

% =============================================================================
\section{Discusión}
\label{sec:discusion}
% =============================================================================

Los resultados del levantamiento de requisitos son coherentes con la evidencia
documentada por Outhwaite~\cite{outhwaite2022apps} y por Topali
et~al.~\cite{topali2025aied}, quienes sostienen que los entornos digitales para
la educación básica deben priorizar interfaces visualmente adaptadas y
mecanismos de soporte que reduzcan la carga cognitiva de los usuarios
infantiles. A diferencia de plataformas como Snappet~\cite{faber2017dfa}, el
modelo adaptativo de Yuhana et~al.~\cite{yuhana2024adaptive} o
MathAIde~\cite{rodrigues2025mathaide}, cuyos requisitos fueron derivados de
contextos y currícula externos al Ecuador, el diseño de Kiddy Quiz partió de un
diagnóstico situado en el sistema educativo ecuatoriano, lo que diferencia su
pertinencia pedagógica de la de los antecedentes internacionales revisados.

La combinación funcional implementada supera en cobertura a las seis
herramientas representativas del estado del arte analizadas en el
Capítulo~II. La revisión sistemática evidenció que ninguna de ellas
---Snappet~\cite{faber2017dfa,molenaar2016snappet}, el modelo de Yuhana
et~al.~\cite{yuhana2024adaptive}, MathAIde~\cite{rodrigues2025mathaide},
\textit{Class Optimization Master}~\cite{zheng2023genai},
MLAD~\cite{mohseni2025mlad} y PictoAndes~\cite{carriontoro2025pictoandes}---
integra simultáneamente la creación docente interna de evaluaciones, la
diversidad de tipos de pregunta, el soporte pictográfico mediante ARASAAC, la
retroalimentación generativa con inteligencia artificial y la gamificación, ni
ha sido validada con SUS y TAM en muestras de estudiantes y docentes de
educación básica elemental latinoamericana.

La puntuación SUS de $87{,}7$ obtenida por el grupo infantil supera la media de
$M=76{,}18$ reportada por Vlachogianni y Tselios~\cite{vlachogianni2022sus}
para sistemas educativos digitales, lo que sugiere que las decisiones de diseño
centradas en pictogramas y simplicidad visual redujeron la carga cognitiva de
los usuarios infantiles en consonancia con los principios de la Teoría Cognitiva
del Aprendizaje Multimedia de Mayer y Moreno~\cite{mayer2003multimedia}. La
brecha registrada en la Facilidad de Uso Percibida entre estudiantes y docentes
no constituye una disfunción del sistema, sino la confirmación empírica de que
la interfaz cumplió su propósito de adaptación al perfil infantil.

La elevada puntuación en Intención de Uso es atribuible a la combinación de
retroalimentación motivacional con elementos lúdicos: mecanismos cuya
efectividad la literatura previa había documentado de manera
independiente~\cite{topali2025aied,mousavinasab2018its,kiili2018}, pero que en
el presente trabajo se integran y validan empíricamente de forma conjunta en un
único instrumento orientado a la educación básica elemental.

\subsection{Contribuciones del estudio}
\label{sec:contribuciones}

\subsubsection*{Contribución teórica}

El estudio contribuye al cuerpo de literatura sobre tecnología educativa en
tres dimensiones. En primer lugar, extiende la evidencia sobre el uso de
pictogramas como soporte cognitivo ---documentada predominantemente en contextos
clínicos y de comunicación aumentativa~\cite{paolieri2018arasaac}--- al ámbito
de la evaluación curricular automatizada en matemáticas para la educación
básica. En segundo lugar, aporta evidencia empírica sobre la viabilidad de
integrar modelos generativos de lenguaje en la retroalimentación formativa
dirigida a niños, dimensión escasamente explorada en la literatura sobre
AIEd K-12~\cite{topali2025aied,mousavinasab2018its}. En tercer lugar,
proporciona datos de validación TAM y SUS con población infantil,
complementando la base empírica revisada por Vlachogianni y
Tselios~\cite{vlachogianni2022sus}, concentrada mayormente en usuarios
adolescentes y adultos.

\subsubsection*{Contribución práctica}

El proyecto entrega un artefacto tecnológico funcional, desplegado y validado,
disponible para instituciones educativas ecuatorianas. La plataforma automatiza
tareas de gestión evaluativa que manualmente resultan inviables para un docente
con numerosos estudiantes ---especialmente la generación de reportes
individualizados, la detección de brechas de aprendizaje por competencia y la
entrega de retroalimentación motivacional inmediata---, respondiendo
directamente a las necesidades manifestadas en la fase de levantamiento de
requisitos y alineándose con los lineamientos del Currículo Priorizado con
Énfasis en Competencias~\cite{mineduc2021curriculo}.

% =============================================================================
\section{Limitaciones del estudio}
\label{sec:limitaciones}
% =============================================================================

La interpretación de los resultados debe considerar las limitaciones inherentes
al alcance del proyecto.

La validación empírica se realizó con una muestra intencional de tamaño
reducido ($N=15$ estudiantes y $N=5$ docentes), seleccionada por accesibilidad
y disposición a participar. Este tamaño es coherente con los estándares mínimos
recomendados para evaluaciones SUS~\cite{bangor2008sus} y con la práctica
habitual en estudios piloto de tecnología educativa, pero no permite realizar
inferencias estadísticas robustas hacia la población general. La generalización
de los hallazgos requiere replicación en muestras mayores, idealmente con
diseños cuasi-experimentales que incorporen grupo control.

El período de exposición de los usuarios al sistema fue acotado. Los efectos
de novedad documentados en la literatura~\cite{vlachogianni2022sus} podrían
haber influido en la elevada puntuación de Intención de Uso. Estudios
longitudinales serán necesarios para confirmar la sostenibilidad de la
aceptación a lo largo del tiempo y el impacto efectivo en los aprendizajes.

La evaluación se circunscribió a la asignatura de matemáticas y al nivel de
educación básica elemental. La extensión a otras asignaturas curriculares y a
otros niveles educativos constituyen extrapolaciones no validadas en el presente
trabajo.

Aunque la integración de Gemini~2.5~Flash demostró ser técnicamente viable, no
se evaluó de manera sistemática la calidad pedagógica de la retroalimentación
generada ni los posibles sesgos del modelo, dimensiones que la literatura
reciente sobre AIEd identifica como áreas críticas~\cite{topali2025aied}.

Finalmente, el desarrollo se ejecutó de manera individual, lo que implicó
adaptar las prácticas colaborativas de XP ---como la programación en pares---
hacia mecanismos de autodisciplina y revisión personal, una limitación que el
propio Beck~\cite{beck2004xp} reconoce como un escenario que reduce
parcialmente las garantías de calidad asociadas a la metodología.

% =============================================================================
\section{Trabajo futuro}
\label{sec:trabajo-futuro}
% =============================================================================

Las limitaciones identificadas y los hallazgos del estudio sugieren cinco
líneas prioritarias de investigación y desarrollo:

\begin{enumerate}
    \item \textbf{Estudios cuasi-experimentales longitudinales.} Diseñar
    estudios con grupo control, aleatorización por aula y mediciones repetidas
    durante al menos un período académico completo, con el propósito de
    cuantificar el efecto de \textit{Kiddy Quiz} sobre indicadores de
    aprendizaje matemático y sobre la disminución de la ansiedad evaluativa,
    controlando el efecto novedad.

    \item \textbf{Extensión curricular y multiidioma.} Adaptar la plataforma a
    otras asignaturas del currículo priorizado~\cite{mineduc2021curriculo}
    ---particularmente lengua y ciencias naturales--- y habilitar
    configuraciones multilingües para contextos plurinacionales como el
    ecuatoriano, incorporando lenguas originarias.

    \item \textbf{Auditoría pedagógica de la retroalimentación generativa.}
    Desarrollar un protocolo de evaluación sistemática de la retroalimentación
    producida por Gemini~2.5~Flash, considerando dimensiones como pertinencia
    pedagógica, ausencia de sesgos, adecuación al nivel cognitivo del estudiante
    y consistencia entre evaluaciones.

    \item \textbf{Personalización adaptativa basada en analítica de
    aprendizaje.} Integrar capacidades de aprendizaje adaptativo que ajusten
    dinámicamente la dificultad y el tipo de pregunta en función del desempeño
    individual del estudiante, en línea con las propuestas de sistemas
    tutoriales inteligentes~\cite{mousavinasab2018its}.

    \item \textbf{Validación con estudiantes con necesidades educativas
    específicas.} Profundizar el componente de accesibilidad mediante
    evaluación con estudiantes con discapacidad intelectual leve, dislexia o
    trastorno del espectro autista, en colaboración con equipos
    interdisciplinarios de psicopedagogía y educación especial.
\end{enumerate}
